\section{Apéndice: índice de términos y ruta de repaso}

\subsection{Términos clave: índice de palabras clave de este episodio}

\noindent\begin{longtable}{p{0.23\textwidth}p{0.32\textwidth}p{0.36\textwidth}}
\toprule
Término & Explicación en una frase & Papel en este episodio \\
\midrule
OpenClaw & Producto y framework de código abierto que detonó el consenso sobre los frameworks de Agent & Representa la exhibición del paradigma de Agent convertido en producto. \\
Post-train & Entrenamiento que continúa después del preentrenamiento y gira en torno a tareas, preferencias, entornos y retroalimentación & Principal terreno de competencia en la era de los Agent. \\
RL Infra & Infraestructura que sostiene el rollout, el reward, el policy update y la eval & Determina si el posentrenamiento de los Agent puede escalar (scale). \\
Rollout & Generación de trayectorias y procesos de interacción en el entorno de la tarea & Fuente de los datos y de la retroalimentación de RL. \\
MiMo-V2 & Familia de modelos de Xiaomi (小米), caso central de la entrevista & Muestra cómo un framework de Agent orquesta varios modelos. \\
Boleto de entrada de 1T & Umbral de escala del modelo base necesario para equipararse con las capacidades de Agent de frontera & Explica por qué el modelo base sigue siendo importante. \\
Igualdad organizacional & Varias funciones se integran en un mismo ciclo de investigación & Explica por qué el posentrenamiento necesita diversity. \\
El entorno importa más que la experiencia & La velocidad de investigación proviene de la capacidad de cómputo, la infra, la organización y el entorno de retroalimentación & Juicio estratégico del cierre. \\
\bottomrule
\end{longtable}

\subsection{Ruta de repaso}

Se recomienda leer primero el capítulo 2 para entender la transición de paradigma de Chat a Agent; después, el capítulo 3 para dominar el ciclo de posentrenamiento, y luego los capítulos 4--6 para entender la orquestación de MiMo, la asignación de capacidad de cómputo y la igualdad organizacional. Los capítulos 7--8 sirven como juicio sobre la industria: por qué en 2026 la competencia ya no se limita al modelo base, sino que consiste en la combinación de arquitectura, posentrenamiento, entorno y capacidad organizacional.

\begin{knowledgebox}{Conexión con EP139}
La entrevista con Su Yu (苏煜) en EP139 ofrece la historia técnica de los Agent y el framework de universal digital agent; la entrevista con Luo Fuli (罗福莉) en EP138 muestra cómo un laboratorio de modelos convierte ese framework en posentrenamiento, orquestación de modelos, asignación de recursos y reestructuración organizacional. Ambos episodios deben leerse juntos.
\end{knowledgebox}

\subsection{Resumen de la sección}

Las palabras clave de este episodio pueden condensarse en seis: OpenClaw, posentrenamiento, RL infra, orquestación de varios modelos, reasignación de capacidad de cómputo e igualdad organizacional. Todas apuntan a un mismo juicio: en la era de los Agent, la competencia es una competencia entre sistemas.

